\section{Introduction}\label{sec:intro}

Logistic regression remains one of the most widely used models for binary
classification, and fitting it is a smooth, strongly convex optimization problem
when an $\ell_2$ penalty is added. In practice the problem is often
\emph{ill-conditioned}: features are strongly correlated, measured on very different
scales, or both, so that the Hessian of the objective has eigenvalues spread over
many orders of magnitude. First-order methods then converge slowly, because their
step size is limited by the largest curvature while progress along low-curvature
directions is proportional to the smallest one~\cite{wright2022optimization,berahas2020investigation}.
Newton-type methods are insensitive to this spread in theory, but they require
second-order information, can take poor steps far from the solution, and, in their
textbook form, solve a dense linear system whose cost grows as $O(nd^2+d^3)$.

A second, more mundane difficulty is numerical. Evaluating
$\log(1+e^{-z})$ and the sigmoid $1/(1+e^{-z})$ directly overflows once $|z|$
exceeds roughly $709$ in double precision and $88$ in single precision. Such
margins arise when features are unscaled, when the starting point is far from the
solution, or when the data are nearly separable, and an overflowing loss silently
breaks line searches and acceptance tests~\cite{blanchard2021accurately,wiemann2024softplus}.

This report studies \emph{Adaptive Stable Hybrid Logistic Regression} (ASH-LR), a
method that combines (i) overflow-free evaluation of the loss and its derivatives,
(ii) an initial phase of inexpensive gradient steps followed by a one-way switch to
second-order steps, and (iii) a damped inexact Newton phase in which the system
$(H+\mu I)p=-g$ is solved by conjugate gradients (CG) using Hessian-vector
products only, and the damping $\mu$ is adapted by a ratio test that compares the
actual with the predicted decrease. In the final version the damping follows the
gradient-regularized rule $\mu_k=\sqrt{\theta_k\norm{g_k}}$, which is motivated by
recent regularized Newton methods~\cite{mishchenko2023regularized,doikov2024super}
and vanishes automatically near the solution.

The project went through two stages. A preliminary version (Milestone~2) used
direct linear solves and line-search-based damping; it converged on all test
problems but needed more iterations and more time than plain Newton. Exploratory
runs then showed that the halving damping rule kept steps over-damped for many
iterations. The final method replaces it, and the final evaluation was fixed in a
pre-registered plan, committed to the code repository before the final experiments
were run, to avoid selecting favourable test cases after the fact.

\paragraph{Research question} Can combining gradient steps with a damped,
matrix-free Newton--CG phase improve the robustness and the work-to-accuracy of
ill-conditioned logistic regression, compared with established matrix-free
methods (Newton--CG and L-BFGS)? The pre-registered success criterion (H1) is a
reduction of at least $20\%$ in the median number of data passes relative to the
best converging matrix-free baseline in at least two of five predeclared
ill-conditioned regimes.

\paragraph{Contributions}
\begin{itemize}
  \item A complete, documented and tested implementation of ASH-LR and five
  baselines with a common interface, an exact work counter in units of data passes,
  and a generator of synthetic designs with an exactly prescribed condition number.
  \item A pre-registered evaluation on eight regimes (six synthetic, two based on the
  Spambase data set) with ten seeds each, complemented by robustness, ablation,
  single-precision and stability experiments.
  \item An honest assessment: we report where ASH-LR helps, where it does not, and
  why direct Newton remains the method of choice for small dense problems.
\end{itemize}

The rest of the report is organized as follows. \Cref{sec:related} reviews related
work, \cref{sec:method} formulates the problem and the algorithm,
\cref{sec:setup} describes the experimental setup, \cref{sec:results} presents and
discusses the results, and \cref{sec:conclusion} concludes.
